\section{The full scalar LDP for stationary equal-rate OU processes}\label{fp:app:ldp}
Let $E,W$ be independent stationary standard OU processes with covariance
$e^{-|s-t|/2}$, and retain $q_t(E)$ and $S_t=-\log q_t(E)$ from
\eqref{fp:eq:basic}. We prove the exact negative-moment domain, regularity
of its pressure, and a full large-deviation principle. The stationary
initial law and equal covariance functions of the two channels are part
of this model.

\begin{theorem}[Full scalar large deviations]\label{fp:thm:scalar-ldp}
Fix $0<\rho<1$, let $\beta^2=\rho/(1-\rho)$ and set
\[
 \theta_c=\beta^{-2}=\frac{1-\rho}{\rho}.
\]
The limit
\begin{equation}\label{fp:eq:Lambda}
 \Lambda_\rho(\theta)
 =\lim_{t\to\infty}\frac1t\log\EE e^{\theta S_t}
 =\lim_{t\to\infty}\frac1t\log\EE q_t^{-\theta}
\end{equation}
exists and is finite for every $\theta<\theta_c$, and
\begin{equation}\label{fp:eq:Lambda-c11}
 \Lambda_\rho\in C^{1,1}_{\mathrm{loc}}(( -\infty,\theta_c)).
\end{equation}
For every $t>0$, $\EE q_t^{-\theta}=\infty$ when $\theta\ge\theta_c$. Furthermore,
\begin{equation}\label{fp:eq:Lambda-boundaries}
 \Lambda_\rho(\theta)\to+\infty,\quad
 \Lambda_\rho'(\theta)\to+\infty\quad(\theta\uparrow\theta_c),
 \qquad
 \Lambda_\rho'(\theta)\to0\quad(\theta\to-\infty).
\end{equation}
The variables $Y_t=S_t/t$ satisfy a full LDP at speed $t$ with good convex rate
\begin{equation}\label{fp:eq:rate}
 I_\rho(y)=\sup_{\theta<\theta_c}\{\theta y-\Lambda_\rho(\theta)\}.
\end{equation}
Explicitly, for every Borel set $A\subset\R$,
\begin{align}
 -\inf_{A^\circ}I_\rho
 &\le\liminf_{t\to\infty}\frac1t\log\PP(Y_t\in A)\notag\\
 &\le\limsup_{t\to\infty}\frac1t\log\PP(Y_t\in A)
 \le-\inf_{\overline A}I_\rho.\label{fp:eq:full-ldp}
\end{align}
\end{theorem}

\subsection{The exact negative-moment domain}
On a finite grid with covariance $\Sigma$, simultaneous whitening gives
\[
 q(x)=\PP\{Z+\beta x\in C\},\qquad
 C=\Sigma^{-1/2}\R_+^n,
\]
where $x,Z$ are standard normal. Since $C$ is a cone,
\[
 q=T_{\sqrt\rho}\ind_C,
 \qquad T_rf(x)=\EE f(rx+\sqrt{1-r^2}Z).
\]
We use the Gaussian reverse-hypercontractive inequality
\cite[Section~3.3, equation~(3.14)]{CDP}:
\begin{equation}\label{fp:eq:reverse-HC}
 \|T_rf\|_s\ge\|f\|_p,
 \qquad s<p<1,\qquad
 r^2\le\frac{1-p}{1-s}.
\end{equation}
Negative exponent norms mean $\|f\|_s=(\EE f^s)^{1/s}$ for positive $f$.
Reverse hypercontractivity originates with Borell~\cite{Borell1982}.
We give the logarithmic Sobolev calculation, following the semigroup
method in \cite{MosselOleszkiewiczSen2013}, to fix the Gaussian
normalization. For the OU semigroup $P_u=T_{e^{-u}}$, set
$r(u)=1-(1-p)e^{2u}$ and
$v_u=P_uf>0$. Then
\begin{align*}
 \frac d{du}\log\|v_u\|_{r(u)}
 ={}&\frac{r'(u)}{r(u)^2}
       \frac{\Ent(v_u^{r(u)})}{\EE v_u^{r(u)}}\\
 &-(r(u)-1)
       \frac{\EE[v_u^{r(u)-2}|\nabla v_u|^2]}{\EE v_u^{r(u)}}.
\end{align*}
Gross's Gaussian logarithmic Sobolev inequality~\cite{Gross1975} gives
\[
 \Ent(v^r)\le\frac{r^2}{2}\EE[v^{r-2}|\nabla v|^2].
\]
Because $r'=2(r-1)<0$, substitution reverses the direction in the first term and makes the derivative nonnegative. Start with bounded strictly positive $f$, pass through exponent zero by the geometric-mean limit, and approximate an indicator from above. This proves the form needed in \eqref{fp:eq:reverse-HC}.

For $0<\theta<\theta_c$, take
\[
 s=-\theta,\qquad p=1-\rho(1+\theta)>0.
\]
Since a negative power reverses the norm inequality, one obtains
\[
 \EE q_{\mathrm{grid}}^{-\theta}
 \le a_{\mathrm{grid}}^{-\theta/[1-\rho(1+\theta)]}.
\]
Along nested dense grids, $q_{\mathrm{grid}}\downarrow q_t$ and $a_{\mathrm{grid}}\downarrow a(t)>0$. Monotone convergence therefore gives
\begin{equation}\label{fp:eq:neg-bound}
 \EE q_t^{-\theta}
 \le a(t)^{-\theta/[1-\rho(1+\theta)]}<\infty.
\end{equation}
In particular,
\begin{equation}\label{fp:eq:Lambda-upper}
 \limsup_t\frac1t\log\EE e^{\theta S_t}
 \le\frac{\theta}{2[1-\rho(1+\theta)]}.
\end{equation}

For the converse, the initial constraint gives $q_t(E)\le\Phi(\beta E_0)$. The Gaussian lower-tail estimate implies that $\EE\Phi(\beta G)^{-\theta}$ dominates a positive constant times
\[
 \int^\infty x^\theta
 \exp\left[\frac{\theta\beta^2-1}{2}x^2\right]dx.
\]
It is infinite when $\theta\beta^2\ge1$, including equality. Thus the finite-time moment domain is exact. Section~\ref{fp:subsec:boundary-divergence} proves the corresponding divergence of the limiting pressure.

\subsection{Covariance interpolation and the existence of all required limits}
Consider two consecutive parts of a finite ordered OU grid. Multiply the one adjacent correlation joining them by $v\in[0,1]$. The new covariance $\Sigma_v$ has unchanged within-block entries and cross-block entries multiplied by $v$. Hence
\[
 \dot\Sigma_{ii}=0,\qquad\dot\Sigma_{ij}\ge0\quad(i\ne j).
\]
Indeed an ordered OU grid has the representation
$X_{k+1}=r_kX_k+\sqrt{1-r_k^2}\,G_{k+1}$, with independent standard
$G_{k+1}$ and $0<r_k<1$. Replacing only the joining $r_k$ by $vr_k$
has exactly the stated effect on the cross-block covariance. The
precision remains tridiagonal with nonpositive off-diagonal entries,
including at $v=0$. Thus its Gaussian density is MTP$_2$; see
\cite[Section~4.1]{Fallat}. Write
\[
 u_v(z)=\PP\{N(0,\Sigma_v)\le z\}.
\]
MTP$_2$ and marginalization give
\begin{equation}\label{fp:eq:cdf-mtp}
 u_i\ge0,\qquad\partial_{ij}\log u\ge0\quad(i\ne j).
\end{equation}
Indeed the set $\{(w,z):w_i\le z_i\ \forall i\}$ is closed under
coordinatewise minimum and maximum. Its indicator times the Gaussian
density is MTP$_2$ in $(w,z)$, and marginalization preserves MTP$_2$
by \cite[Proposition~3.4]{Fallat}. An independent product density in
$z$ can be inserted and cancelled to obtain an integrable joint density.

For $\eta>-\beta^{-2}$ set
\[
 H_\eta(v)=\EE_{X_v\sim N(0,\Sigma_v)}u_v(\beta X_v)^\eta.
\]
Differentiate the environmental covariance using Gaussian integration by parts, and the private covariance using its heat equation. The two contributions are
\begin{equation}\label{fp:eq:interpolation}
 H_\eta'(v)
 =\frac\eta2\sum_{i,j}\dot\Sigma_{ij}
 \EE\left[u^{\eta-2}
 \{(1+\beta^2)u u_{ij}+\beta^2(\eta-1)u_i u_j\}\right],
\end{equation}
where derivatives of $u$ are evaluated at $\beta X_v$. The bracket rewrites exactly as
\begin{equation}\label{fp:eq:interpolation-positive}
 (1+\beta^2)u^2\partial_{ij}\log u
 +(1+\eta\beta^2)u_i u_j.
\end{equation}
Only $i\ne j$ terms occur. By \eqref{fp:eq:cdf-mtp} and $1+\eta\beta^2>0$, the bracket is nonnegative. Thus
\begin{equation}\label{fp:eq:interpolation-sign}
 H_\eta'(v)\ge0\quad(\eta>0),\qquad
 H_\eta'(v)\le0\quad(-\beta^{-2}<\eta<0).
\end{equation}
To justify differentiation for negative powers, first fix the grid.
The covariance path is uniformly positive definite. After whitening,
Lemma~\ref{fp:lem:curvature} bounds first and second spatial derivatives
of $\log u$ by fixed polynomials, uniformly in $v$. The
reverse-hypercontractive proof above applies to every $\Sigma_v$.
Its ordinary orthant probability is bounded below uniformly in $v$:
Gaussian association gives at least $2^{-n}$ for a grid of $n$ points.
Choose $\theta'$ with $-\eta<\theta'<\beta^{-2}$. The resulting
uniform $\theta'$ inverse-moment bound and H\"older's inequality
control every polynomial-weighted differentiated integrand.
One may first differentiate with smooth cutoffs and then remove them
using this uniform integrability. The positive-power case follows by
the same polynomial bounds and $u^\eta\le1$.

At $v=0$ the two blocks are independent, and at $v=1$ their original
joint law is restored. For adjacent continuous intervals choose nested
dense grids with no duplicated boundary point. In each block, points
approach the common endpoint from its own side; continuity then imposes
that endpoint constraint in both marginal limits. For negative powers,
monotone convergence and the finite bound \eqref{fp:eq:neg-bound}
justify this limit. If $M_\theta(t)=\EE e^{\theta S_t}$, we obtain
\begin{equation}\label{fp:eq:theta-subadd}
 M_\theta(s+t)\le M_\theta(s)M_\theta(t),
 \qquad0<\theta<\theta_c,
\end{equation}
whereas
\begin{equation}\label{fp:eq:theta-superadd}
 M_\theta(s+t)\ge M_\theta(s)M_\theta(t),
 \qquad\theta<0.
\end{equation}
For positive $\theta$, \eqref{fp:eq:neg-bound} bounds $\log M_\theta(t)$ above by $C(t+1)$ and it is nonnegative. For negative $\theta$, \eqref{fp:eq:linear-b} bounds it below by $-C(t+1)$ and it is nonpositive. Continuous-parameter subadditivity or superadditivity therefore proves existence and finiteness in \eqref{fp:eq:Lambda}; more precisely,
\begin{equation}\label{fp:eq:Lambda-inf-sup}
 \Lambda_\rho(\theta)=
 \begin{cases}
 \inf_{t>0}t^{-1}\log M_\theta(t),&0<\theta<\theta_c,\\
 0,&\theta=0,\\
 \sup_{t>0}t^{-1}\log M_\theta(t),&\theta<0.
 \end{cases}
\end{equation}

\subsection{Differentiability throughout the finite domain}
On a whitened grid, write $A=-\log q$ and
\[
 d\nu_\theta\propto e^{\theta A}d\gamma,
 \qquad V_\theta(x)=|x|^2/2-\theta A(x).
\]
Lemma~\ref{fp:lem:curvature} gives
\begin{equation}\label{fp:eq:kappa-theta}
 D^2V_\theta\succeq\kappa_\theta I,\qquad
 \kappa_\theta=
 \begin{cases}
 1,&\theta\le0,\\
 1-\theta\beta^2,&0<\theta<\theta_c.
 \end{cases}
\end{equation}
For $\theta\le0$, convexity of $A$ gives the lower bound $I_d$. Brascamp--Lieb and \eqref{fp:eq:curvature} imply
\[
 \Var_{\nu_\theta}(A)
 \le\frac{2\beta^2}{\kappa_\theta}\EE_{\nu_\theta}A.
\]
Consequently, for $\Lambda_{t,\mathrm{grid}}=t^{-1}\log\EE e^{\theta A}$,
\begin{equation}\label{fp:eq:Lambda-curvature}
 0\le\Lambda_{t,\mathrm{grid}}''(\theta)
 \le\frac{2\beta^2}{\kappa_\theta}
       \Lambda_{t,\mathrm{grid}}'(\theta).
\end{equation}
The moment bounds on a slightly larger compact subinterval of $(-\infty,\theta_c)$ and convex secants give a uniform bound for $\Lambda_{t,\mathrm{grid}}'$ for $t\ge1$. Thus the second derivatives are uniformly bounded on compacts. A neighboring exponent dominates $A^je^{\theta A}$, $j\le2$, when passing to dense grids. The same bounds hold for the continuous-time $\Lambda_t$.

Apply Arzel\`a--Ascoli to the first derivatives and use
\[
 \Lambda_t(b)-\Lambda_t(a)=\int_a^b\Lambda_t'(\theta)\,d\theta
\]
together with pointwise convergence. Every subsequential derivative limit is the derivative of $\Lambda_\rho$. Hence \eqref{fp:eq:Lambda-c11} holds and
\begin{equation}\label{fp:eq:Lambda-derivative-limit}
 \Lambda_t'\longrightarrow\Lambda_\rho'
 \quad\text{locally uniformly on }(-\infty,\theta_c).
\end{equation}

\subsection{Extensive divergence at the finite boundary}
\label{fp:subsec:boundary-divergence}
Take a unit-spaced grid of $n$ times in $[0,t]$, with $n/t\to1$. Its covariance is $\Sigma_{ij}=r^{|i-j|}$, $r=e^{-1/2}$. Let $Q=\Sigma^{-1}$. Both matrices have eigenvalues in $[K^{-1},K]$, with $K=(1+r)/(1-r)$; this also follows from the explicit tridiagonal precision and row-sum bounds.

For an environmental vector satisfying $Qx\le0$, choose the nonnegative Chernoff vector $h=-\beta Qx$. Then
\begin{equation}\label{fp:eq:chernoff-cone}
 q_{\mathrm{grid}}(x)=\PP\{Z\ge-\beta x\}
 \le e^{\beta h^Tx+\frac12h^T\Sigma h}
 =e^{-\beta^2x^TQx/2}.
\end{equation}
Since $q_t\le q_{\mathrm{grid}}$, for $0<\theta<\theta_c$ and $\delta=1-\theta\beta^2$,
\begin{align}
 M_\theta(t)
 &\ge\int_{\{Qx\le0\}}e^{\theta\beta^2x^TQx/2}\,dN(0,\Sigma)(x)\notag\\
 &=\delta^{-n/2}\PP\{QX\le0\},\qquad X\sim N(0,\Sigma).\label{fp:eq:boundary-integral}
\end{align}
The equality uses homogeneity of the cone. Now $QX\sim N(0,Q)$; on $[-1,0]^n$ its density is at least
\[
 (2\pi K)^{-n/2}e^{-Kn/2}.
\]
Thus $\PP(QX\le0)\ge e^{-Cn}$, and
\begin{equation}\label{fp:eq:boundary-divergence}
 \Lambda_\rho(\theta)
 \ge\frac12\log\frac1{1-\theta\beta^2}-C.
\end{equation}
This proves divergence of $\Lambda_\rho$ at $\theta_c$; convexity proves divergence of its derivative. The linear number of grid constraints makes the divergence extensive in time.

For the other boundary, fix $a>0$ and set
\[
 p_a=\PP(E_s\ge a\ (0\le s\le1)),\qquad
 r_a=\PP(W_s\ge-\beta a\ (0\le s\le1)).
\]
Both are positive. Association on consecutive unit intervals gives
\[
 \EE q_t^\eta\ge p_a^{\lceil t\rceil}r_a^{\eta\lceil t\rceil},
 \qquad \Lambda_\rho(-\eta)\ge\log p_a+\eta\log r_a.
\]
First let $\eta\to\infty$, then $a\to\infty$, using $r_a\to1$. Since $\Lambda_\rho(-\eta)\le0$, this yields $\Lambda_\rho(-\eta)/(-\eta)\to0$. Convexity proves $\Lambda_\rho'(-\infty)=0$. This completes \eqref{fp:eq:Lambda-boundaries}.

\subsection{The full upper and lower LDP bounds}
We now apply the G\"artner--Ellis argument~\cite{DemboZeitouni2010}
to the pressure just constructed, giving the upper and lower bounds
explicitly.

For any fixed $0<\theta_0<\theta_c$, Chernoff gives
\[
 \limsup_{t\to\infty}\frac1t\log\PP(Y_t>R)
 \le\Lambda_\rho(\theta_0)-\theta_0R\longrightarrow-\infty.
\]
Since $Y_t\ge0$, the laws are exponentially tight.

By continuity of $\Lambda_\rho'$ and its two boundary limits, each $y>0$ is $\Lambda_\rho'(\theta)$ for some interior parameter. Convexity then gives
\[
 I_\rho(y)=\theta y-\Lambda_\rho(\theta).
\]
Under
\[
 \frac{d\PP_{\theta,t}}{d\PP}
 =e^{\theta S_t-t\Lambda_t(\theta)},
\]
the mean of $Y_t$ tends to $y$ by \eqref{fp:eq:Lambda-derivative-limit}, and its variance is $\Lambda_t''(\theta)/t=O(t^{-1})$. Thus $\PP_{\theta,t}(|Y_t-y|<\epsilon)\to1$. Changing measure back yields
\[
 \liminf_{t\to\infty}\frac1t\log\PP(|Y_t-y|<\epsilon)
 \ge-I_\rho(y)-|\theta|\epsilon.
\]
Let $\epsilon\downarrow0$ to obtain the local lower bound. The endpoint $y=0$ is obtained by positive approximation if $I_\rho(0)<\infty$: convexity and lower semicontinuity then imply right continuity there. If $I_\rho(0)=\infty$ it adds no lower-bound obligation. For $y<0$, \eqref{fp:eq:rate} is infinite, by the sublinear left asymptotic of $\Lambda_\rho$.

For the upper bound, each fixed $\theta$ gives a Chernoff bound on a sufficiently small neighborhood of a point. Choose an almost-optimal supporting parameter in \eqref{fp:eq:rate}, cover a compact set by finitely many such neighborhoods, and use exponential tightness outside that compact set. This proves the closed-set upper bound. The rate is convex and lower semicontinuous; it dominates $\theta_0y-\Lambda_\rho(\theta_0)$ on $y\ge0$, so its sublevel sets are compact. Theorem~\ref{fp:thm:scalar-ldp} is proved.

\subsection{Consequences for the spectrum and tilted environments}
Define the extended spectrum
\[
 \widetilde\chi(\eta,\rho)=-\Lambda_\rho(-\eta),
 \qquad \eta>-(1-\rho)/\rho.
\]
It agrees with \eqref{fp:eq:positive-spectrum} on positive orders. Convex duality gives
\begin{equation}\label{fp:eq:duality-extended}
 \widetilde\chi(\eta,\rho)=\inf_{y\ge0}\{I_\rho(y)+\eta y\},\qquad
 I_\rho(y)=\sup_{\eta>-(1-\rho)/\rho}
 \{\widetilde\chi(\eta,\rho)-\eta y\}.
\end{equation}
The minimizer in the first formula is uniquely $\Lambda_\rho'(-\eta)$: equality in convex duality forces $y$ into the singleton subgradient of the differentiable function $\Lambda_\rho$ at $-\eta$. This proves \eqref{fp:eq:scalar-var}.

The unique zero of $I_\rho$ is $\gamma_\rho=\Lambda_\rho'(0)$, and
\begin{equation}\label{fp:eq:quenched-derivative}
 \partial_\eta\widetilde\chi(0,\rho)=\gamma_\rho
 =\lim_{t\to\infty}\frac1t\EE[-\log q_t].
\end{equation}
Small positive and negative Chernoff parameters give exponential concentration about $\gamma_\rho$. Borel--Cantelli along integers and monotonicity of $S_t$ imply $S_t/t\to\gamma_\rho$ almost surely. Jensen and \eqref{fp:eq:Lambda-upper} give
\[
 \frac12\le\gamma_\rho\le\frac1{2(1-\rho)}.
\]
Theorem~\ref{fp:thm:annealing} improves the lower bound to
\begin{equation}\label{fp:eq:gamma-strict}
 \gamma_\rho\ge\frac12+\frac\rho{16\pi},
\end{equation}
since for $0<\eta\le1/2$ it bounds $\chi(\eta,\rho)/\eta$ from below, and this quotient tends to $\partial_\eta\widetilde\chi(0,\rho)$.

For the pure $q_t^\eta$-tilted OU environment, $Y_t$ has the full good rate
\begin{equation}\label{fp:eq:tilted-rate}
 I_{\eta,\rho}(y)=I_\rho(y)+\eta y-\chi(\eta,\rho),\qquad\eta>0.
\end{equation}
One proves this by the same change of measure on compact sets; the nonnegative linear tilt and exponential tightness control the upper tail. Its unique zero is $\partial_\eta\chi(\eta,\rho)$. The mixed arithmetic source has the separate local conclusion in
Theorem~\ref{fp:thm:history-local-ldp}.

Finally,
\begin{equation}\label{fp:eq:tail-slope}
 \lim_{y\to\infty}\frac{I_\rho(y)}y
 =\theta_c=\frac{1-\rho}{\rho}.
\end{equation}
Every $\theta<\theta_c$ gives the lower limiting slope. For $y\ge\gamma_\rho$, negative $\theta$ cannot improve the dual supremum over its value at zero, because $\Lambda_\rho(\theta)\ge\gamma_\rho\theta$. For $\theta\ge0$, $\Lambda_\rho(\theta)\ge0$, so $I_\rho(y)\le\theta_cy$. This proves the opposite inequality.

The exact Lamperti representation $E_u=e^{-u/2}B^E_{e^u}$, $W_u=e^{-u/2}B^W_{e^u}$ transfers the scalar statement to Brownian random-wall survival on $[1,R]$ with speed $\log R$. The initial law is the distribution induced by Brownian time $1$.
